% Einheit 1 von 3 – Potenz-, Faktor- und Summenregel (bank/ableitungsregeln/e1.jsonl, zusammenbau v0.1)
%% TODO Zeitmarke Einheit 1: Marken-Zeile nicht lesbar
%% TODO Zweigzeile Teil 1: was hier gelernt wird (Satz in Schülersprache aus dem Einheitstitel) – Einheit 1
\einheitenkopf[e1]{Einheit 1 von 3 · Potenz-, Faktor- und Summenregel}
\zweigzeile{Abitur GK}
\setcounter{aufgabe}{9}

%% TODO Titel: Ich-kann-Satz fehlt in der Bank (Platzhalter: Kettenname) – Nr. 10
%% TODO Anweisung: ein Satz über der ersten Teilaufgabe fehlt in der Bank – Nr. 10
\begin{aufgabe}{Zahl oder Variable?}
\begin{teile}
\teil Für jede Zahl $a > 0$ ist $f_a(x) = a x^4 - 3x$. Welcher Buchstabe ist beim Ableiten nach $x$ eine feste Zahl? \\ \kreuz{$a$} \\ \kreuz{$x$} \\ \kreuz{$a$ und $x$}
\end{teile}
\end{aufgabe}

%% TODO Titel: Ich-kann-Satz fehlt in der Bank (Platzhalter: Kettenname) – Nr. 11
%% TODO Anweisung: ein Satz über der ersten Teilaufgabe fehlt in der Bank – Nr. 11
\begin{aufgabe}{Potenz-, Faktor- und Summenregel}
%% TODO \rechenplatz{n}: Schrittzahl des Grundfalls fehlt in der Bank (nur bei fester Schreibform, 2.3 b)
\begin{teile}
\teil $f(x) = x^9$ – welche Regel reicht? \\ \kreuz{Potenzregel} \\ \kreuz{Faktorregel} \\ \kreuz{Konstantenregel}
\teil $f(x) = 2x^3 + 5x^2 - 4x + 9$ – $f'(x)$? $f'(x) =$ \leerfeld
\teil $f(x) = \frac{1}{6}x^3 + \frac{3}{2}x^2 - 2x$ – $f'(x)$? $f'(x) =$ \leerfeld
\teil $f(x) = -2x^5 + 3{,}25x^4 + 6x^3 - 30x^2$ – erste, zweite und dritte Ableitung? \hfill (FHR 2026) \\ $f'(x) =$ \leerfeld , $f''(x) =$ \leerfeld , $f'''(x) =$ \leerfeld
\teil $f(x) = -3x^5 + 4x^3 - x$ – Ableitung $f'(x)$ und eine Stammfunktion $F(x)$? \hfill (Abitur 2019 GK) \\ $f'(x) =$ \leerfeld , $F(x) =$ \leerfeld
\teil $f(x) = \frac{1}{3}x^3 - \frac{9}{2}x + 4$ – die ersten drei Ableitungen und der Anstieg des Graphen im Schnittpunkt mit der $y$-Achse? \hfill (FHR 2024) \\ $f'(x) =$ \leerfeld , $f''(x) =$ \leerfeld , $f'''(x) =$ \leerfeld , Anstieg = \leerfeld
\teil Für jede Zahl $a > 0$ ist $f_a(x) = x^4 - a \cdot x^3$. Zeige, dass $f_a'(x) = x^2 \cdot (4x - 3a)$ gilt. \hfill (Abitur 2021 LK)
\teil $f(x) = x^3 - 9x^2 + 24x$. Weise nach: $f'(x) = 3 \cdot (x - 2) \cdot (x - 4)$. \hfill (Abitur 2025 LK)
\teil $f(x) = \frac{1}{5}x^5 - 6x^3 + 81x$. Weise nach: $f'(x) = (x - 3)^2 \cdot (x + 3)^2$. \hfill (Abitur 2022 GK)
\teil Gegeben ist $f(x) = \frac{1}{4}x^4 + x^3 - 2x^2$. Bestimme die ersten drei Ableitungen. Prüfe, ob $P(2|4)$ auf dem Graphen von $f$ liegt und ob $P$ ein Extrempunkt ist. \hfill (FHR 2024) \\ $f'(x) =$ \leerfeld , $f''(x) =$ \leerfeld , $f'''(x) =$ \leerfeld
\end{teile}
\end{aufgabe}

%% TODO Titel: Ich-kann-Satz fehlt in der Bank (Platzhalter: Kettenname) – Nr. 12
%% TODO Anweisung: ein Satz über der ersten Teilaufgabe fehlt in der Bank – Nr. 12
\begin{aufgabe}{Verhalten im Unendlichen bestimmen}
\begin{teile}
\teil $f(x) = 2x^3 + x^2 - 12x - 5$ – Verhalten der Funktionswerte für $x \to \pm\infty$ und die ersten drei Ableitungen? \hfill (FHR 2022) \\ $f'(x) =$ \leerfeld , $f''(x) =$ \leerfeld , $f'''(x) =$ \leerfeld
\end{teile}
\end{aufgabe}

%% TODO Titel: Ich-kann-Satz fehlt in der Bank (Platzhalter: Kettenname) – Nr. 13
%% TODO Anweisung: ein Satz über der ersten Teilaufgabe fehlt in der Bank – Nr. 13
\begin{aufgabe}{Potenz-, Faktor- und Summenregel – Fehler finden, Begründen}
\begin{teile}
\teil Ich finde den Fehler: Zu $f(x) = 3x^4 - 2x^2 + 6$ schreibt Mia $f'(x) = 12x^4 - 4x^2$.
\teil Begründe, warum die Ableitung einer Konstanten, etwa $f(x) = 12$, null ist.
\end{teile}
\end{aufgabe}
